% ============================================================================
% 第二章：方法
% ============================================================================

\section{方法概述}

\begin{frame}{基于表格型 Q-Learning 的时间窗优化方法}
    \begin{block}{核心思想}
        将时间窗优化问题建模为马尔可夫决策过程 (MDP)，使用表格型 Q-Learning 算法学习最优时间窗选择策略
    \end{block}
    \vspace{0.3cm}
    \begin{itemize}
        \item \textbf{状态空间}: 136 个离散化的时间窗参数对
        \item \textbf{动作空间}: 4 个时间窗边界调整操作
        \item \textbf{奖励函数}: CSP-SVM 分类器的交叉验证准确率
        \item \textbf{状态转移}: 确定性过程
    \end{itemize}
\end{frame}

% ============================================================================
% 问题形式化
% ============================================================================
\section{问题形式化}

\begin{frame}{马尔可夫决策过程建模}
    \begin{block}{状态空间 $\mathcal{S}$}
        \[
        \mathcal{S} = \left\{ (t_{\text{start}}, t_{\text{end}}) \mid t_{\text{end}} - t_{\text{start}} \geq 0.5 \right\}
        \]
        \begin{itemize}
            \item 离散化步长：$\Delta t = 0.2$ 秒
            \item 有效状态数量：\textbf{136 个}
        \end{itemize}
    \end{block}
    
    \vspace{0.3cm}
    \begin{block}{动作空间 $\mathcal{A}$}
        \[
        \mathcal{A} = \{a_0, a_1, a_2, a_3\}
        \]
        \begin{itemize}
            \item $a_0$: 左边界左移 (变宽)
            \item $a_1$: 左边界右移 (变窄)
            \item $a_2$: 右边界左移 (变窄)
            \item $a_3$: 右边界右移 (变宽)
        \end{itemize}
    \end{block}
\end{frame}

\begin{frame}{奖励函数与状态转移}
    \begin{block}{奖励函数 $\mathcal{R}$}
        \[
        r(s, a) = \text{Accuracy}_{\text{CV}}(s'; \mathcal{D}_{\text{train}})
        \]
        \begin{itemize}
            \item 使用 5 折交叉验证计算准确率
            \item 基础分类器：CSP-SVM
        \end{itemize}
    \end{block}
    
    \vspace{0.3cm}
    \begin{block}{状态转移 $\mathcal{P}$}
        \[
        \mathcal{P}(s' \mid s, a) = 
        \begin{cases}
            1, & \text{if } s' = f(s, a) \\
            0, & \text{otherwise}
        \end{cases}
        \]
        \begin{itemize}
            \item 确定性转移 (离散化带来的便利)
        \end{itemize}
    \end{block}
\end{frame}

% ============================================================================
% 算法设计
% ============================================================================
\section{算法设计}

\begin{frame}{表格型 Q-Learning 算法}
    \begin{block}{Q 值更新规则}
        \[
        Q(s_t, a_t) \leftarrow Q(s_t, a_t) + \alpha \left[ r_{t+1} + \gamma \max_{a} Q(s_{t+1}, a) - Q(s_t, a_t) \right]
        \]
        \begin{itemize}
            \item 学习率 $\alpha = 0.2$
            \item 折扣因子 $\gamma = 0.95$
            \item TD 误差：当前估计值与目标值的差距
        \end{itemize}
    \end{block}
\end{frame}

\begin{frame}{探索与利用策略}
    \begin{block}{$\varepsilon$-贪婪策略}
        \begin{itemize}
            \item 以概率 $1-\varepsilon$ 选择当前 Q 值最大的动作 (利用)
            \item 以概率 $\varepsilon$ 随机选择动作 (探索)
        \end{itemize}
    \end{block}
    
    \vspace{0.3cm}
    \begin{exampleblock}{参数设置}
        \begin{itemize}
            \item 初始探索率：$\varepsilon_{\text{init}} = 0.9$
            \item 最小探索率：$\varepsilon_{\text{min}} = 0.05$
            \item 衰减率：$\varepsilon_{\text{decay}} = 0.995$
            \item 训练轮数：100 episodes (4000 步)
        \end{itemize}
    \end{exampleblock}
\end{frame}

\begin{frame}{算法流程}
    \begin{itemize}
        \item 初始化 Q 表：$Q[s, a] \leftarrow 0$
        \item 对于每个 episode (1--100):
        \begin{itemize}
            \item 随机初始化状态 $s_0$
            \item 对于每个 step (1--40):
            \begin{itemize}
                \item 根据$\varepsilon$-贪婪选择动作 $a_t$
                \item 执行动作，状态转移 $s_{t+1} \leftarrow f(s_t, a_t)$
                \item 计算奖励 $r_t$ (5 折交叉验证)
                \item 更新 Q 值
            \end{itemize}
            \item 衰减探索率 $\varepsilon$
        \end{itemize}
        \item 返回最优状态 $s^* = \arg\max_s \max_a Q[s, a]$
    \end{itemize}
\end{frame}

\begin{frame}{推理过程}
    \begin{block}{最优时间窗提取}
        \[
        s^* = \arg\max_{s \in \mathcal{S}} \max_{a \in \mathcal{A}} Q(s, a)
        \]
    \end{block}
    
    \vspace{0.3cm}
    \begin{exampleblock}{计算复杂度}
        \begin{itemize}
            \item 时间复杂度：$O(|\mathcal{S}| \cdot |\mathcal{A}|) = O(544)$
            \item 实际计算中可忽略不计
            \item 满足在线 BCI 系统实时性要求
        \end{itemize}
    \end{exampleblock}
\end{frame}

% ============================================================================
% 收敛性分析
% ============================================================================
\section{收敛性分析}

\begin{frame}{收敛性定理}
    \begin{block}{Q-Learning 收敛性 (Watkins \& Dayan, 1992)}
        对于有限 MDP，若满足：
        \begin{enumerate}
            \item 所有状态 - 动作对被无限次访问
            \item 学习率序列满足 Robbins-Monro 条件
            \item 采用适当的探索策略
        \end{enumerate}
        则 Q 值序列以概率 1 收敛到 $Q^*$
    \end{block}
    
    \vspace{0.3cm}
    \begin{alertblock}{本研究的收敛性保证}
        \begin{itemize}
            \item 有限状态空间：136 个状态 × 4 个动作 = 544 条目
            \item 100 轮 × 40 步 = 4000 步，平均每对访问约 7.4 次
            \item $\varepsilon$-贪婪策略保证充分探索
        \end{itemize}
    \end{alertblock}
\end{frame}

% ============================================================================
% 方法对比
% ============================================================================
\section{方法对比}

\begin{frame}{表格型 Q-Learning vs 深度 Q 网络 (DQN)}
    \begin{table}[h]
        \centering
        \footnotesize
        \begin{tabular}{l|c|c}
            \hline
            \textbf{特性} & \textbf{表格型 Q-Learning} & \textbf{DQN} \\
            \hline
            参数量 & 544 (Q 表) & $\sim 10^6$ (神经网络) \\
            收敛性保证 & \checkmark 理论保证 & $\times$ 无保证 \\
            可解释性 & 完全可解释 & 黑盒模型 \\
            训练稳定性 & 稳定 & 对超参数敏感 \\
            样本效率 & 高 & 低 \\
            计算资源 & 低 & 高 \\
            \hline
        \end{tabular}
        \caption{表格型 Q-Learning 与 DQN 对比}
    \end{table}
\end{frame}

\begin{frame}{适用边界}
    \begin{block}{表格型 Q-Learning 适用场景}
        \begin{itemize}
            \item 状态空间规模 $< 10^4$
            \item 状态 - 动作对可充分访问
            \item 需要可解释性和稳定性
        \end{itemize}
    \end{block}
    
    \vspace{0.3cm}
    \begin{block}{DQN 适用场景}
        \begin{itemize}
            \item 高维连续状态空间
            \item 状态空间过大无法表格化
            \item 有充足数据和计算资源
        \end{itemize}
    \end{block}
\end{frame}
